\subsection{FREE MONOIDS}

Any finite sequence \(w = (x_{i})_{1 \leqslant i \leqslant n}\) of elements of \(\mathbf{X}\) indexed by an interval \([1, n]\) of \(\mathbf{N}\) (possibly empty) is called a word constructed on \(\mathbf{X}\) . The integer \(n\) is called the length of the word \(w\) and denoted by \(l(w)\) . There is a unique word of length 0, namely the empty sequence e. X will be identified with the set of words of length 1.

Let \(w = (x_{i})_{1 \leqslant i \leqslant m}\) and \(w' = (x_{j}')_{1 \leqslant j \leqslant n}\) be two words. The composition of \(w\) and \(w'\) is the word \(u = (y_{k})_{1 \leqslant k \leqslant m + n}\) defined by

\[
y _ {k} = \left\{ \begin{array}{l l} x _ {k} & \text { for } 1 \leqslant k \leqslant m \\ x _ {k - m} ^ {\prime} & \text { for } m + 1 \leqslant k \leqslant m + n. \end{array} \right.\tag{1}
\]

In other words, the sequence \(w''\) is obtained by first writing the elements of the sequence w and then those of \(w'\) . The composition of w and \(w'\) is generally denoted by \(ww'\) or \(w.w'\) ; it is sometimes said that it is obtained by juxtaposition of w and \(w'\) . Then by construction \(l(w.w') = l(w) + l(w')\) .

The relation we = ew = w is immediately established for every word w. Let \(w = (x_{i})_{1 \leqslant i \leqslant m}\) , \(w' = (x_{j}')_{1 \leqslant j \leqslant n}\) and \(w'' = (x_{k}'')_{1 \leqslant k \leqslant p}\) be three words; clearly the words \(w(w'w'')\) and \((ww')w''\) are both equal to the word \((y_{i})_{1 \leqslant i \leqslant m+n+p}\) defined by

\[
y _ {l} = \left\{ \begin{array}{l l} x _ {l} & \text { if } 1 \leqslant l \leqslant m \\ x _ {l - m} ^ {\prime} & \text { if } m + 1 \leqslant l \leqslant m + n \\ x _ {l - m - n} ^ {\prime \prime} & \text { if } m + n + 1 \leqslant l \leqslant m + n + p. \end{array} \right.\tag{2}
\]

The above shows that the set of words constructed on X with the law of composition \((w, w') \mapsto w.w'\) is a monoid with identity element e. It is denoted by \(\mathrm{Mo}(\mathbf{X})\) and called the free monoid constructed on X. It follows immediately from the definition of product of words that every word \(w = (x_i)_{1 \leq i \leq n}\) is equal to the product \(\prod_{i=1}^{n} x_i\) . A word may therefore be written in the form \(x_1 \ldots x_{n}\) .

PROPOSITION 3. Let M be a monoid. Every mapping f of X into M extends uniquely to a homomorphism of Mo(X) into M.

Let \(g\) be a homomorphism of \(\operatorname{Mo}(\mathbf{X})\) into \(\mathbf{M}\) extending \(f\) . If \(w = (x_{i})_{1 \leqslant i \leqslant n}\) is a word, then \(w = \prod_{i=1}^{n} x_{i}\) in the monoid \(\operatorname{Mo}(\mathbf{X})\) , whence

\[
g (w) = \prod_ {i = 1} ^ {n} g (x _ {i}) = \prod_ {i = 1} ^ {n} f (x _ {i})
\]

in the monoid M (§ 1, no. 2, formula (2)). This proves the uniqueness of \(g\) .

Let \(h(w) = \prod_{i=1}^{n} f(x_i)\) for every word \(w = (x_i)_{1 \leqslant i \leqslant n}\) . The associativity theorem (§ 1, no. 3, Theorem 1) and the definition of product in Mo(X) imply \(h(ww') = h(w)h(w')\) . By convention the empty product \(h(e)\) is the identity element of M and \(h(x) = f(x)\) for \(x \in X\) . Hence \(h\) is a homomorphism of Mo(x) into M extending \(f\) .

Let \(u: X \to Y\) be a mapping. By Proposition 3, there exists one and only one homomorphism of \(\operatorname{Mo}(\mathbf{X})\) into \(\operatorname{Mo}(\mathbf{Y})\) which coincides with \(u\) on \(\mathbf{X}\) ; it is denoted by \(\operatorname{Mo}(u)\) . It maps a word \((x_i)_{1 \leqslant i \leqslant n}\) to the word \((u(x_i))_{1 \leqslant i \leqslant n}\) . As in the case of magmas (no. 1), the equation \(\operatorname{Mo}(v \circ u) = \operatorname{Mo}(v) \circ \operatorname{Mo}(u)\) is established for every mapping \(v: \mathbf{Y} \to \mathbf{Z}\) and it can be shown that \(\operatorname{Mo}(u)\) is injective (resp. surjective, bijective) if \(u\) is. For every subset S of \(\mathbf{X}\) , \(\operatorname{Mo}(\mathbf{S})\) is identified with the submonoid of \(\operatorname{Mo}(\mathbf{X})\) generated by S.

Let X be a set and \((u_{\alpha}, v_{\alpha})_{\alpha \in I}\) be a family of ordered pairs of elements of \(\mathrm{Mo}(\mathrm{X})\) . Let R be the equivalence relation on \(\mathrm{Mo}(\mathrm{X})\) compatible with the law on \(\mathrm{Mo}(\mathrm{X})\) and generated by the \((u_{\alpha}, v_{\alpha})\) (§1, no. 6). The monoid \(\mathrm{Mo}(\mathrm{X})/\mathrm{R}\) is called the monoid defined by X and the relators \((u_{\alpha}, v_{\alpha})_{\alpha \in I}\) . Let h be the canonical morphism of \(\mathrm{Mo}(\mathrm{X})\) onto \(\mathrm{Mo}(\mathrm{X})/\mathrm{R}\) . Then \(\mathrm{Mo}(\mathrm{X})\) is generated by \(h(\mathrm{X})\) .

Let N be a monoid and \((n_{x})_{x\in\mathbf{X}}\) a family of elements of N. Let k be the morphism of \(\mathrm{Mo}(\mathbf{X})\) into N such that \(k(x)=n_{x}\) for all \(x\in X\) (Proposition 3). If \(k(u_{\alpha})=k(v_{\alpha})\) for all \(\alpha\in I\) , there exists one and only one magma morphism \(f:\mathrm{Mo}(\mathbf{X})/\mathbf{R}\to\mathbf{N}\) such that \(f(h(x))=n_{x}\) for all \(x\in X\) (§1, no. 6, Proposition 9); as k is unital, f is a monoid morphism.
% semantic-fragments: legacy-input-seam
\relax{}
